\section{Introduction}
\label{sec:introduction}

An ideal is said to be integrally closed if it coincides with its integral
closure, whereas it is said to be normal if all its positive powers are
integrally closed.  Thus an integrally closed ideal need not be normal.  Since
the integral closedness of an ideal alone provides no general procedure for
controlling its higher powers, it is natural to ask what restrictions on the
number and form of the generators ensure that an integrally closed ideal
is normal.  When the ambient ring is a normal domain, the normalization of the
Rees algebra of $I$ is $\bigoplus_{n\geq0}\ol{I^n}T^n$.  Hence normality of
$I$ is equivalent to normality of its Rees algebra, so this question is also
naturally connected with the study of normal blowup algebras.

Results organized by the number of generators provide a natural point of
comparison.  Let $(R,\m)$ be a regular local ring of dimension $d$, let
$I$ be an integrally closed $\m$-primary ideal, and write $\mu_R(I)$ for the
minimal number of generators of $I$.  Goto \cite{Goto1987} proved that the Rees
algebra of $I$ is a Cohen--Macaulay normal domain when $\mu_R(I)=d$, that is,
when $I$ is a complete intersection
\cite[Corollary~1.3]{Goto1987}.  Ciuperc\u{a} \cite{Ciuperca2006} obtained the same conclusion when
$\mu_R(I)=d+1$
\cite[Theorem~1.1]{Ciuperca2006}.  Endo, Goto, Hong, and Ulrich \cite{EndoGotoHongUlrich2026} extended these
results by proving the conclusion whenever $\mu_R(I)\leq d+2$
\cite[Theorem~2.2 and Corollary~2.4]{EndoGotoHongUlrich2026}.  They further proved that, over a field of characteristic zero, the Rees algebra of an
integrally closed zero-dimensional ideal generated by $d+3$ homogeneous
polynomials in the standard graded polynomial ring $k[x_1,x_2,\ldots,x_d]$ is a
Cohen--Macaulay normal domain
\cite[Theorem~3.4]{EndoGotoHongUlrich2026}.  Thus, in three variables, their
results cover every integrally closed zero-dimensional ideal with at most five
minimal generators, and also cover six-generated ideals in the standard
graded case in characteristic zero.

For monomial ideals in three variables, Ataka and the author
\cite{AtakaMatsuoka2026} proved that, over an arbitrary field, every
integrally closed zero-dimensional monomial ideal with at most seven minimal
generators is normal \cite[Theorem~3.1]{AtakaMatsuoka2026}.  
This bound is sharp: the integral closure \(\overline{(x^7,y^3,z^2)}\) is 
minimally generated by eight monomials but is not normal; 
see \cite[Exercise 1.14]{SwansonHuneke2006} and \cite[Example 4.5]{AtakaMatsuoka2026}. 
We revisit this ideal in Example \ref{ex:successive-cutoffs-eight-generators}, where it occurs as the next cutoff after a normal ideal.
Together with the standard facts recalled in the next paragraph, the theorem of
Ataka--Matsuoka also implies that the corresponding Rees algebra is a
Cohen--Macaulay normal domain.  Consequently, among six-generated ideals in
three variables, the first case not covered by the results above is the
non-monomial case outside the standard graded setting; in positive characteristic, even the
standard graded nonmonomial case lies outside the theorem of
Endo--Goto--Hong--Ulrich.

The monomial case has two particularly powerful features.  First, a theorem
of Reid, Roberts, and Vitulli states that a monomial ideal in $d$ variables is
normal once its first $d-1$ positive powers are integrally closed
\cite[Proposition~3.1]{ReidRobertsVitulli2003}.  In particular, for an
integrally closed monomial ideal in three variables, it is enough to verify
the integral closedness of its square.  Second, the Rees algebra of a monomial
ideal is an affine semigroup ring.  Once the ideal is normal, its Rees algebra
is therefore a normal affine semigroup ring and is Cohen--Macaulay by
Hochster's theorem \cite{Hochster1972}.  Neither result applies directly to an
ideal containing a binomial generator.  In particular, the finite criterion
of Reid--Roberts--Vitulli does not reduce the problem studied here to the
integral closedness of the square.  Moreover, when the relevant grading is not
standard, the theorem of Endo--Goto--Hong--Ulrich does not apply either.  Thus,
outside both the monomial and standard graded settings, the preceding results
provide neither a finite criterion for normality nor a direct conclusion that
the Rees algebra is normal.

These results suggest the following question.

\begin{question}\label{que:seven-generators}
Let $k$ be an arbitrary field, let $S=k[x,y,z]$, and let $I$ be an integrally
closed $(x,y,z)$-primary ideal that is homogeneous with respect to some
positive weighted grading of $S$.  If $\mu_S(I)\leq7$, must $I$ be normal?
\end{question}

The purpose of this paper is to study a natural class of six-generated
integrally closed ideals that lies just beyond these monomial and standard
graded settings.  Recall that a numerical semigroup is a submonoid
$H$ of $\mathbb N$ such that $\mathbb N\setminus H$ is finite. Here, $\mathbb{N}$ denotes the set of all non-negative integers.
Its multiplicity $\mult(H)$ is its least positive element, and its embedding
dimension is the cardinality of its minimal system of generators.  Writing
$\Fr(H)=\max(\mathbb N\setminus H)$, called the Frobenius number of $H$, we say that $H$ is symmetric if
$a \in H$ if and only if $\Fr(H)-a \notin H$ for every integer $a$.  We
determine the ideals arising from the numerical semigroup rings $k[H]$ when
$H$ has embedding dimension three and multiplicity three or four.  In the
multiplicity-three and symmetric multiplicity-four cases, apart from the
cases obtained by adjoining one variable to an ideal in a two-variable
polynomial ring, the resulting six-generated ideals form two explicit
families.  We give an affirmative answer to
Question~\ref{que:seven-generators} for these two families by proving that
their Rees algebras are Cohen--Macaulay normal domains.

We first state the main result directly in terms of the resulting ideals.
Their numerical-semigroup origin is explained after Theorem~\ref{thm:main}.  Let $k$ be
an arbitrary field, let $S=k[x,y,z]$, and put $\m=(x,y,z)$.  Consider either
of the following ideals:
\begin{equation*}\tag{$\sharp$}\label{eq:family-one}
\begin{aligned}
 I={}&(y^2-x^{b-d}z,\ x^{2b+2-e},\ x^{b+1}y,\
        x^{b-d+1}z,\ yz,\ z^2),\\
 &\hspace{35mm} b>d\geq0,\qquad e\in\{0,1\},
\end{aligned}
\end{equation*}
or
\begin{equation*}\tag{$\sharp\sharp$}\label{eq:family-two}
\begin{aligned}
 I={}&(y^2-x^{2b+1},\ x^{2b+2},\ x^{b+1}y,\ x^rz,\ yz,\ z^2),\\
 &\hspace{35mm} b\geq1,\qquad 1\leq r\leq b+1.
\end{aligned}
\end{equation*}
Each of these ideals is homogeneous with respect to a suitable positive
grading.  For example, one may take $\deg x=3$, $\deg y=3b+2-e$, and
$\deg z=3b+3d+4-2e$ for \eqref{eq:family-one}, while one may take
$\deg x=4$, $\deg y=4b+2$, and any positive value of $\deg z$ for
\eqref{eq:family-two}.  Each ideal is $\m$-primary and minimally generated
by six elements.  By
Theorem~\ref{thm:classification}(4), every ideal of either form arises from
the numerical-semigroup construction explained after Theorem~\ref{thm:main}
and is therefore integrally closed.

\begin{theoremA}\label{thm:main}
For every ideal $I$ of either form \eqref{eq:family-one} or
\eqref{eq:family-two}, the Rees algebra of $I$ is a Cohen--Macaulay normal
domain.
\end{theoremA}

The ideal in \eqref{eq:family-one} is homogeneous for the standard grading
precisely when $b-d=1$, whereas no ideal in \eqref{eq:family-two} is
homogeneous for the standard grading.  When $b-d=1$, the
characteristic-zero case is covered by the theorem of
Endo--Goto--Hong--Ulrich, but the positive-characteristic case is not.  The
ideals in \eqref{eq:family-one} with $b-d>1$, as well as all the ideals in
\eqref{eq:family-two}, lie outside the scope of that theorem in every
characteristic.  Thus Theorem~\ref{thm:main} goes beyond the ranges covered by the two
results above, including in positive characteristic along the standard graded
boundary.

We now explain the origin of these ideals.  Let
$H=\langle a_1,a_2,a_3\rangle$ be a numerical semigroup of embedding
dimension three, where $a_1,a_2,a_3$ form its minimal system of generators,
let $k[H]=k[t^{a_1},t^{a_2},t^{a_3}]$ be the $k$-subalgebra of the polynomial ring $k[t]$, and consider the graded homomorphism
\[
 \varphi_H\colon S=k[x_1,x_2,x_3]\longrightarrow k[H],
 \qquad x_i\longmapsto t^{a_i},
\]
where $\deg x_i=a_i$.  For each $0<h\in\mathbb N$, set
\[
 I_h=\varphi_H^{-1}\bigl(t^hk[t]\cap k[H]\bigr).
\]
The ideal $t^hk[t]$ is integrally closed in the polynomial ring $k[t]$.
Hence its contraction to $k[H]$, and therefore its inverse image
$I_h$, is integrally closed.

Let $\fkp_H=\Ker\varphi_H$, and let $\ell$ be the least weighted degree of a
nonzero homogeneous element of $\fkp_H$.  If $h\leq\ell$, then $I_h$ is a
monomial ideal.  The descriptions of the defining ideals of numerical
semigroup rings of embedding dimension three used here are based on Herzog's work
\cite{Herzog1970}.  For the numerical semigroups considered in this paper,
the degree-$\ell$ component of $\fkp_H$ is one-dimensional over $k$ and is
spanned by a binomial.  Consequently, when $h=\ell+1$, precisely this relation
remains as a binomial generator, while the relations of larger degree are
absorbed into the monomial part of $I_{\ell+1}$.  This is why the ideals in
Theorem~\ref{thm:main} have five monomial generators and one binomial generator.
The present paper focuses on those $I_{\ell+1}$ that are minimally generated
by six elements, and investigates whether the relations among their
generators provide enough control over arbitrary powers to prove normality.

The following classification shows that the two families in Theorem~\ref{thm:main}
are precisely the six-generated families arising from this construction.
When comparing $I_{\ell+1}$ with the ideals
in \eqref{eq:family-one} and \eqref{eq:family-two}, we use $x,y,z$ for a
suitable ordering of $x_1,x_2,x_3$.

\begin{theoremB}\label{thm:classification}
Let $H$ be a numerical semigroup of embedding dimension three and multiplicity
three or four, let $\ell$ be as above, and consider $I_{\ell+1}$.
\begin{enumerate}[label=\textup{(\arabic*)}]
 \item If $H$ has multiplicity three, then, after a possible interchange of
 variables, $I_{\ell+1}$ is of the form \eqref{eq:family-one}.
 \item If $H$ is symmetric of multiplicity four, equivalently, if $k[H]$ is
 Gorenstein, then, after a possible interchange of variables, either
 \[
  I_{\ell+1}=(x_3)+JS
 \]
 for an integrally closed $(x_1,x_2)$-primary ideal
 $J\subseteq k[x_1,x_2]$, or $I_{\ell+1}$ is of the form
 \eqref{eq:family-one} or \eqref{eq:family-two}.
 \item If $H$ is non-symmetric of multiplicity four, then $I_{\ell+1}$ is
 minimally generated by seven elements.
 \item Conversely, every choice of the parameters in
 \eqref{eq:family-one} is realized by a numerical semigroup of multiplicity
 three, and every choice of the parameters in \eqref{eq:family-two} is
 realized by a symmetric numerical semigroup of multiplicity four.
\end{enumerate}
\end{theoremB}

Thus Theorem~\ref{thm:main} covers every six-generated ideal arising from this construction
when the semigroup has multiplicity three or is symmetric of multiplicity
four, apart from the cases that reduce immediately to two variables.  We
emphasize that Theorem~\ref{thm:classification} classifies the ideals produced by this
numerical-semigroup construction; it is not a classification of all
six-generated integrally closed ideals that are homogeneous for some positive
grading.

These results provide evidence for the expectation that, in a three-dimensional regular local ring $(A,\mathfrak n)$, 
every six-generated integrally closed $\mathfrak n$-primary ideal has a Cohen–Macaulay normal Rees algebra. 
The two families treated here establish this conclusion for a class of non-monomial ideals over an arbitrary field, including ideals that are not homogeneous for the standard grading.

For Cohen–Macaulayness, we explicitly choose a three-generated ideal $Q\subseteq I$ and verify that $I^2=QI+\mathfrak mI^2$. 
After localization, this equality allows us to apply the standard criterion for Cohen–Macaulayness of Rees algebras. 
The verification uses relations among the generators specific to each family.


The first case of Theorem~\ref{thm:classification}(2) is covered by
\cite[Theorem~2.1]{EndoGotoHongUlrich2026}, as explained in
Subsection~\ref{subsec:symmetric-multiplicity-four}.

The non-symmetric case of multiplicity four lies just beyond the scope of
Theorem~\ref{thm:main} because it produces seven generators.  To give a seven-generated
affirmative instance of Question~\ref{que:seven-generators}, we study in
detail the ideal obtained from the numerical semigroup
$H=\langle4,9,15\rangle$.  After determining its seven minimal generators in
the proof of Theorem~\ref{thm:classification}, we prove that it is normal,
that its localization at $\m$ has reduction number one, and that its Rees
algebra is a Cohen--Macaulay normal domain.  This example provides evidence
for the first part of Question~\ref{que:seven-from-semigroups}, concerning
all seven-generated ideals arising from non-symmetric numerical semigroups of
embedding dimension three and multiplicity four.  Although we do not claim a general theorem for this class,
we expect the same method to apply throughout the non-symmetric
multiplicity-four case. Since the additional case analysis is unlikely to
yield further mathematical insight, we restrict the detailed treatment to
one seven-generated example.

The construction also suggests a question about successive degree bounds.
For $h\in H$, set
\[
 \sigma(h)=\min\{n\in H\mid n>h\}.
\]
Assume that both $I_h$ and $I_{\sigma(h)}$ have at most seven minimal
generators.  If $I_h$ is normal, must $I_{\sigma(h)}$ also be normal?
Without the generator bounds, this implication can fail even for two
successive monomial ideals, as Example~\ref{ex:successive-cutoffs-eight-generators}
shows.  The successor
$\sigma(h)$, rather than $h+1$, is the natural quantity here because $h+1$
need not belong to $H$, and several consecutive integers may therefore define
the same ideal.  Although $\ell+1$ need not belong to $H$, one always has
$I_{\ell+1}=I_{\sigma(\ell)}$.  The passage from $I_\ell$ to
$I_{\ell+1}$ is the first passage from a monomial ideal to
an ideal with one binomial generator.  
Because the strategy of Subsection \ref{subsec:normality} uses the normality of $J=I+(f)$ independently of how it is established, it may be applied successively to the pairs $I_{\sigma(h)}\subset I_h$, provided that the product reductions and coefficient-separation conditions can be verified at each step.

The paper is organized as follows.  Section~\ref{sec:general-setup} fixes the
numerical-semigroup setup and explains the two common proof strategies: the
reduction calculation used to prove the Cohen--Macaulayness of the Rees
algebra and the coefficient-separation argument used to prove normality.  In
Section~\ref{sec:target-ideals} we determine the ideals $I_{\ell+1}$ arising
in multiplicities three and four and prove
Theorem~\ref{thm:classification}.
Section~\ref{sec:main-theorem} proves
Theorem~\ref{thm:main}; the families in \eqref{eq:family-one} and
\eqref{eq:family-two} are treated in Subsections~\ref{sec:family-one}
and~\ref{sec:family-two}, respectively.
Section~\ref{sec:nonsymmetric-example} is devoted to the seven-generated ideal
arising from $\langle4,9,15\rangle$.
Section~\ref{sec:further-questions} records several questions suggested by
these results, including the question concerning successive ideals $I_h$.
